%%% FIELD primitive-family-correction
For \(\theta=(\mu,f,b,x,y,\zeta,r,t_1,t_2)\in\Theta_{\mathcal G}^{\mathrm{leg}}\) and \(d\in\mathbb R\), put
\[
 \bar b=b+\zeta,\qquad B=\mu^2+f,\qquad
 R=t_1+B+x,\qquad \eta=B/R,\qquad K=\mu/B,
\]
\[
 c_1=t_1^{-1/2},\qquad
 t_{2,n}=t_2-\frac{R^2}{B^2}\frac d{\sqrt n},
 \qquad c_{2,n}=t_{2,n}^{-1/2}.
\]
Generate mutually independent coordinates
\[
 F_t\sim N(\mu,f),\quad e_{0t}\sim N(0,b),\quad
 e_{1t}\sim N(0,x),\quad e_{2t}\sim N(0,y),\quad
 \xi_t\sim N(0,\zeta),
\]
independently over \(t\), and define
\[
 Y_t(0)=F_t+e_{0t}+\xi_t,\qquad
 \Delta_t=\tau,\qquad
 W_{1t}=c_1(F_t+e_{1t}),\qquad
 W_{2t}=c_{2,n}(F_t+e_{2t}).
\]
Role one uses \(W_{1t}\) as donor and \(W_{2t}\) as supplemental proxy, while role two exchanges those assignments.  Thus, conditional on \(F_t\), the treated idiosyncratic error \(e_{0t}+\xi_t\) and the two factor-measurement errors are mutually independent, and the causal target satisfies \(\mathbb E_n\Delta_t=\tau\).

The tied manifold is \(t_2=t_1+x-y>0\).  Along the displayed local calibration, \(t_{2,n}>0\) eventually, both bridge matrices equal \(c_1c_{2,n}B\), the bridge coefficients are \(1/c_1\) and \(1/c_{2,n}\), and
\[
 \sqrt n(L_1-L_2)=d+O(n^{-1/2}).
\]
At the tied limit,
\[
 \sigma_D^2
 =4\eta^2\{(1-\eta)^2B+\bar b\}(x+y)
  +2\eta^4(x^2+y^2),
\]
\[
 c_{D1}=\operatorname{Cov}(D,A_1)
 =-2\eta K\{(1-\eta)Bx+\bar b y\},
\qquad
 c_{D2}=\operatorname{Cov}(D,A_2)
 =2\eta K\{(1-\eta)By+\bar b x\},
\]
and, because treatment is constant,
\[
 V_1=K^2(B+y)(\bar b+x)+r(\bar b+x),
\qquad
 V_2=K^2(B+x)(\bar b+y)+r(\bar b+y).
\]
Consequently
\[
 \Gamma(\theta,d)
 =\left(
 -\frac d{\sigma_D},
 \frac{c_{D1}}{\sqrt{\sigma_D^2V_1}},
 \frac{c_{D2}}{\sqrt{\sigma_D^2V_2}},
 V_1,V_2
 \right).
\]
The corrected family \(\mathcal G\) is the collection of these constant-treatment local Gaussian arrays.
%%% FIELD anchor-statement
The corrected Gaussian family in \(\texttt{def:primitive-gaussian-family}\), restricted to any compact neighborhood on which its bridge and fixed-role variance inequalities are strict, is a stationary independent-over-time, one-factor, square proximal synthetic-control subclass satisfying Shi et al.'s Assumptions 1--5, supplementary conditions D.1--D.10, and the fixed-role full-rank condition for either predeclared role.  The same conclusion holds row by row for a diagonal constant-effect family in which the donor and supplemental-proxy measurement errors are independent conditional on the factor and the exceptional centered rare-jump variable occurs only in the treated unit's idiosyncratic untreated-outcome error.  Hence the reduced-form class \(\mathcal M(D_{\mathrm{gap}})\) is not asserted to contain, or to be contained in, the full published class; rather, these explicit laws belong to their intersection.  Consequently, any coverage or impossibility property proved separately for these laws also holds within every published class containing this subclass.
%%% FIELD anchor-proof
We verify the published-model conditions directly.  First consider a fixed row of the corrected family.  Write
\[
 \varepsilon_{0t}=e_{0t}+\xi_t,\qquad
 \varepsilon_{1t}=c_1e_{1t},\qquad
 \varepsilon_{2t}=c_{2,n}e_{2t}.
\]
The observed untreated outcomes have the one-factor representation
\[
 Y_t(0)=1\cdot F_t+\varepsilon_{0t},\qquad
 W_{1t}=c_1F_t+\varepsilon_{1t},\qquad
 W_{2t}=c_{2,n}F_t+\varepsilon_{2t}.
\]
The three errors are mutually independent conditional on \(F_t\), because all five primitive Gaussian coordinates in the corrected construction are mutually independent.  This proves the required donor--proxy conditional independence and the corresponding exclusion condition for either exchange of the two control roles.  Consistency holds by construction: the observed treated outcome is \(Y_t(0)\) before treatment and \(Y_t(0)+\tau\) afterward, with no anticipation.  The treatment effect is the constant \(\Delta_t=\tau\), so the causal estimand is not introduced through an observed moment:
\[
 \mathbb E_n\Delta_t=\mathbb E_n\tau=\tau.
\]

For role \(j\), the scalar synthetic-control loading is \(\alpha_j=1/c_j\), because
\[
 c_j\alpha_j=1.
\]
For the opposite role \(k\), independence and \(B=\mathbb E_nF_t^2>0\) give
\[
 M_s=\mathbb E_n(W_{kt}W_{jt})
     =c_kc_jB>0,
\qquad
 \mathbb E_n\{W_{kt}Y_t(0)\}=c_kB=M_s\alpha_j.
\]
Thus the bridge exists, is unique, and has full row rank.  Its residual is
\[
 Y_t(0)-W_{jt}\alpha_j=\varepsilon_{0t}-e_{jt},
\]
which is centered.  Adding the constant treatment effect and taking expectations yields
\[
 \mathbb E_n\{Y_t(0)+\Delta_t-W_{jt}\alpha_j\}
 =\tau+\mathbb E_n(\varepsilon_{0t}-e_{jt})=\tau.
\]
This proves, rather than defines, the relation between the causal target and the published fixed-role identifying moment.

All observable coordinates are independent and identically distributed over \(t\).  They are therefore strictly stationary and ergodic, their absolute-regularity coefficients vanish after lag zero, and their pre- and post-treatment residual processes have absolutely summable covariance sequences.  Gaussian laws are dominated by Lebesgue measure and have moments of every finite order.  On a compact parameter neighborhood bounded away from
\[
 f=0,\quad b=0,\quad x=0,\quad y=0,\quad \zeta=0,\quad
 t_1=0,\quad t_2=0,\quad M_s=0,
\]
the Gaussian moments possess one common integrable envelope.  The moment maps are polynomials in the data and continuously differentiable in the finite-dimensional bridge and ATT parameters.  Their population bridge derivative is
\[
 -\mathbb E_n(W_{kt}W_{jt})=-M_s\neq0,
\]
and their ATT derivative is \(-1\).  The resulting fixed-role Jacobian is triangular with nonzero diagonal and hence nonsingular.  The law of large numbers gives sample-Jacobian convergence.

For a fixed role, the pre-period bridge score is \(W_{kt}(\varepsilon_{0t}-e_{jt})\) and the post-period ATT residual is \(\varepsilon_{0t}-e_{jt}\).  Their respective variances are strictly positive because
\[
 \mathbb E_n(W_{kt}^2)>0,\qquad
 \operatorname{Var}(\varepsilon_{0t}-e_{jt})
 =b+\zeta+\operatorname{Var}(e_{jt})>0,
\]
and the two time blocks are disjoint.  Hence the fixed-role limiting moment covariance is positive definite even though the covariance obtained by stacking both roles may be singular.  The assumed limit \(n/m\to r\in(0,\infty)\) supplies comparable horizons.  Compactness of a sufficiently small closed neighborhood of the true finite-dimensional parameter, together with the strict Jacobian, variance, and rank inequalities just verified, supplies the remaining compactness, identification, domination, derivative, covariance, and full-rank requirements in supplementary conditions D.1--D.10.

The diagonal family in the maximal-failure theorem has the same one-factor equations.  In an ordinary row its treated idiosyncratic error is Gaussian.  In row \(q\), it is
\[
 \varepsilon_{0t}^{(q)}
 =I_{q,t}-p_q+a_qZ_{0t},
 \qquad p_q=q^{-2},\qquad a_q=q^{-4},
\]
where \(I_{q,t}\) is Bernoulli with success probability \(p_q\), and \(Z_{0t}\) is independent standard normal.  This error is centered, has a strictly positive Lebesgue density because \(a_q>0\), has moments of every finite order, and is independent of the factor and of both Gaussian measurement errors.  The two measurement-error variances in that row are \(a_q^2>0\).  Thus every fixed row retains domination, conditional proxy independence, positive fixed-role covariance, and nonsingular Jacobian and bridge rank.  All rows are independent and identically distributed over time, treatment remains \(\Delta_t=\tau\), and the horizons are comparable.  The same calculations therefore verify the published assumptions row by row for the diagonal family.

These checks establish an explicit common subclass.  They do not prove a set inclusion between the paper's whole reduced-form class and Shi et al.'s full class: the former imposes geometric beta mixing and permits a singular covariance after stacking roles, whereas the latter contains structural interactive-effects and fixed-role conditions not imposed throughout \(\mathcal M(D_{\mathrm{gap}})\).  Any result established separately on the displayed common subclass transfers upward to a class containing it, and the corrected primitive Gaussian points inherit the same published-class interpretation.
%%% FIELD maximal-theorem-statement
The asserted whole-array uniform selector-conditional convergence is false under \(\mathcal M(D_{\mathrm{gap}})\), and its failure can be maximal even inside the stationary independent-over-time square subclass satisfying Shi et al.'s Assumptions 1--5, supplementary conditions D.1--D.10, and the fixed-role full-rank condition for both predeclared roles.  Specifically, fix
\[
 0<\kappa<\frac{6373}{15369},\qquad
 0<\nu<\frac{39}{625},
\]
and a sufficiently large finite moment envelope.  For every sufficiently large integer \(q\), there is a whole triangular array \((P_j^{(q)})_{j\geq1}\in\mathcal M(0)\) with constant treatment effect \(\Delta_t=\tau\), conditionally independent factor-measurement errors, and
\[
 \operatorname{Var}(D)=\frac{39}{625},\qquad
 V_s^{\mathrm{gap}}=V_s=1\quad(s\in\mathcal S),
\]
whose stacked limiting covariance is singular, but for
\[
 T_q=
 \frac{\widehat\tau_{\mathrm{gap}}-\tau}
      {\widehat{\operatorname{se}}_{\widehat J_A}}
\]
one has
\[
 \lim_{q\to\infty}
 \mathbb E_{P_q^{(q)}}\!\left[
  \sup_{x\in\mathbb R}
  \left|
   \Pr_{P_q^{(q)}}\{T_q\leq x\mid\mathcal F_A\}-\Phi(x)
  \right|
 \right]=1.
\]
Thus positive limiting marginal variances, coordinatewise studentization, a logarithm-squared mixing gap, and singular-covariance-safe Gaussian mapping do not yield whole-array uniform inference without a uniform finite-row variance and distributional-stabilization condition.
%%% FIELD maximal-theorem-proof
We construct and verify the arrays.  Let
\[
 B_0=\frac{6373}{30000},\qquad
 x_0=\frac18,\qquad
 \bar b_0=\frac78,\qquad
 t_0=\frac{5123}{10000},\qquad
 c_0=t_0^{-1/2}.
\]
These constants were chosen so that
\[
 t_0+B_0+x_0=4B_0,\qquad
 \eta_0=\frac{B_0}{t_0+B_0+x_0}=\frac14,
\qquad
 \frac{B_0}{t_0}=\frac{6373}{15369}>\kappa.
\]
Set \(m_j=j\), so \(j/m_j=1\) and the sample-ratio condition holds with \(r=1\).

For each sufficiently large \(q\), define a whole array \((P_j^{(q)})_{j\geq1}\).  Every row is independent and identically distributed over the two-sided time index.  In every ordinary row \(j\neq q\), generate mutually independent
\[
 F_t\sim N(0,B_0),\qquad
 e_{1t},e_{2t}\sim N(0,x_0),\qquad
 \varepsilon_{0t}\sim N(0,\bar b_0),
\]
and put
\[
 Y_t(0)=F_t+\varepsilon_{0t},\qquad
 \Delta_t=\tau,\qquad
 W_{1t}=c_0(F_t+e_{1t}),\qquad
 W_{2t}=c_0(F_t+e_{2t}).
\]
Role one uses \(W_1\) as donor and \(W_2\) as proxy; role two exchanges them.  The Gaussian treated error can, if desired, be split as the sum of two independent \(N(0,7/16)\) coordinates, exactly matching the corrected primitive parameterization.

In the exceptional row \(j=q\), retain \(F_t\sim N(0,B_0)\), let
\[
 p_q=q^{-2},\qquad a_q=q^{-4},
\]
and generate mutually independent
\[
 I_{q,t}\sim\operatorname{Bernoulli}(p_q),\qquad
 Z_{0t},Z_{1t},Z_{2t}\sim N(0,1).
\]
Define
\[
 \varepsilon_{0t}=I_{q,t}-p_q+a_qZ_{0t},\qquad
 e_{1t}=a_qZ_{1t},\qquad e_{2t}=a_qZ_{2t},
\]
and use the same formulas for \(Y_t(0)\), \(\Delta_t\), and \(W_{jt}\).  Thus the exceptional centered skewness lies in the untreated treated-unit error, never in the treatment effect.  Conditional on \(F_t\), the treated error and the two measurement errors are mutually independent.  The causal target is
\[
 \mathbb E_{P_j^{(q)}}\Delta_t=\tau
\]
in every row, while the observed post-treatment residual mean equals this target by the bridge calculation below.

We verify membership in \(\mathcal M(0)\).  Independence over time implies strict stationarity and beta-mixing coefficients equal to zero at every positive lag.  The Gaussian coordinates have fixed moments, \(0\leq I_{q,t}\leq1\), and \(p_q,a_q\leq1\); consequently, for every fixed \(\delta_0>0\),
\[
 \sup_q\sup_j
 \mathbb E_{P_j^{(q)}}\|X_t\|^{8+\delta_0}<\infty.
\]
Choose \(C_{\mathrm{mom}}\) larger than this finite bound.

For either role and every row, independence and centering give
\[
 M_s=\mathbb E(Z_{s,t}W_{D_s,t})=c_0^2B_0
     =\frac{6373}{15369}>\kappa.
\]
Moreover,
\[
 \mathbb E\{Z_{s,t}Y_t(0)\}=c_0B_0=M_s/c_0.
\]
Since this scalar moment is nonzero, the bridge definition gives
\[
 \alpha_s=\frac1{c_0},
\qquad
 Y_t(0)-W_{D_s,t}\alpha_s=\varepsilon_{0t}-e_{s,t},
\]
where \(e_{s,t}\) denotes the selected donor's measurement error.  The last residual is centered, proving bridge centering.  With constant treatment,
\[
 \mathbb E\{Y_t(0)+\Delta_t-W_{D_s,t}\alpha_s\}
 =\tau,
\]
so the observed identifying moment equals the separately defined causal target.

Both donor variables have the same distribution within every row.  Hence their population ridge coefficients and minimized population losses agree, so
\[
 L_{1,j}=L_{2,j},\qquad
 \sup_{j\geq1}\sqrt j\,|L_{1,j}-L_{2,j}|=0.
\]
This verifies the bounded-gap condition with \(D_{\mathrm{gap}}=0\).

It remains to calculate the limiting variances.  Since each donor mean is zero, \(\mu_s=0\), so \(H_s=0\), \(\psi_{s,t}=0\), and \(A_s=0\) for both roles.  In every ordinary row, the tied scalar formula from the corrected primitive construction applies with
\[
 B=B_0,\qquad x=y=x_0,\qquad
 \bar b=\bar b_0,\qquad \eta=\eta_0.
\]
Therefore
\[
 \begin{aligned}
 \operatorname{Var}(D)
 &=4\eta_0^2\{(1-\eta_0)^2B_0+\bar b_0\}(2x_0)
   +4\eta_0^4x_0^2\\
 &=\frac{39}{625}.
 \end{aligned}
\]
For transparency, the exact cancellation is
\[
 (1-\eta_0)^2B_0+\bar b_0
 =\frac{159119}{160000},
\]
and substitution of \(B_0=6373/30000\), \(x_0=1/8\), and \(\eta_0=1/4\) gives the displayed rational value.  The post residuals in an ordinary row are
\[
 v_{1,t}=\varepsilon_{0t}-e_{1t},\qquad
 v_{2,t}=\varepsilon_{0t}-e_{2t},
\]
so
\[
 \operatorname{Var}(R_s)=\bar b_0+x_0=1,
\qquad
 V_s=\operatorname{Var}(A_s)+r\operatorname{Var}(R_s)=1>\nu.
\]
By \(\texttt{lem:gap-variance-inheritance}\), \(V_s^{\mathrm{gap}}=V_s=1\), because \(\operatorname{Var}(A_s)=0\).  The stacked limiting covariance is singular because both pre-ATT coordinates vanish identically.

For each fixed \(q\), all rows after \(q\) have the same ordinary law.  Thus the exceptional row cannot alter the array's limiting long-run covariance.  The centered score vector in an ordinary row is a degree-two Gaussian polynomial and has a finite third moment.  For any coefficient vector \(a\), if \(G_t\) denotes that centered score vector and \(\Sigma_G\) its covariance, Taylor expansion gives
\[
 \mathbb E\exp\!\left\{\frac{\mathrm i a^\prime G_t}{\sqrt j}\right\}
 =1-\frac{a^\prime\Sigma_Ga}{2j}+O(j^{-3/2}).
\]
Independence over \(t\) and raising this expression to the \(j\)-th power yield
\[
 \exp\!\left(-\frac12a^\prime\Sigma_Ga\right).
\]
This Cramér--Wold calculation proves covariance convergence and the possibly singular joint Gaussian limit without any covariance inversion.  We have verified every defining condition of \(\mathcal M(0)\).  The rowwise published-class inclusion follows from \(\texttt{prop:shi-anchor-transfer}\): the factor errors are conditionally independent, the bridge is scalar and full rank, treatment is constant, and the exceptional error has a positive Gaussian-convolution density and finite moments.

We now analyze row \(q\).  Let \(\mathcal N_q\) be the event that \(I_{q,t}=0\) for every \(t\) in the late bridge-fitting block \(\mathcal B_q\) and post block \(\mathcal C_q\).  Since their total length is at most \(2q\),
\[
 \Pr(\mathcal N_q)\geq(1-p_q)^{2q}
 =(1-q^{-2})^{2q}\longrightarrow1.
\]
This event is independent of the early sigma-field \(\mathcal F_A\).  The selected label is \(\mathcal F_A\)-measurable, and the following bounds hold for both labels.

Let \(j\) denote the donor chosen by a fixed role and \(k\neq j\) its proxy.  On \(\mathcal N_q\), the true bridge residual on both retained blocks is
\[
 r_{j,t}^0=\varepsilon_{0t}-e_{j,t}
 =-p_q+a_q(Z_{0t}-Z_{jt}).
\]
Writing empirical averages over \(\mathcal B_q\) with a bar and
\[
 \widehat M_{j,q}
 =|\mathcal B_q|^{-1}\sum_{t\in\mathcal B_q}W_{kt}W_{jt},
\]
the sample bridge equation gives the exact identity
\[
 \widehat\alpha_j-\alpha_j
 =\widehat M_{j,q}^{-1}
   |\mathcal B_q|^{-1}
   \sum_{t\in\mathcal B_q}W_{kt}r_{j,t}^0.
\]
Because \(|\mathcal B_q|/q\to1-\varrho>0\),
\[
 \widehat M_{j,q}=c_0^2B_0+O_P(q^{-1/2}),
\qquad
 \overline W_{k,\mathcal B_q}=O_P(q^{-1/2}).
\]
Conditional on \(\mathcal N_q\), the centered Gaussian part of
\(W_{kt}r_{j,t}^0\) has standard deviation \(O(a_q)\), while its constant part is \(-p_qW_{kt}\).  Hence
\[
 \widehat\alpha_j-\alpha_j
 =O_P\!\left(\frac{p_q+a_q}{\sqrt q}\right)
\]
uniformly over the two roles.  The post donor mean satisfies
\[
 \overline W_{j,\mathcal C_q}=O_P(q^{-1/2}).
\]
Consequently the definition of the gap-split estimator yields, on \(\mathcal N_q\),
\[
 \begin{aligned}
 \widehat\tau_{\mathrm{gap}}-\tau
 &=\overline r^0_{j,\mathcal C_q}
   -(\widehat\alpha_j-\alpha_j)
     \overline W_{j,\mathcal C_q}\\
 &=-p_q+O_P(a_q q^{-1/2})
   +O_P\{(p_q+a_q)q^{-1}\}\\
 &=-p_q\{1+o_P(1)\}.
 \end{aligned}
\]

We next bound the actual Bartlett denominator.  For any centered scalar observations \(g_1,\ldots,g_N\), the Bartlett estimator with bandwidth \(h\) satisfies, by Cauchy--Schwarz and weights bounded by one,
\[
 0\leq\widehat\Omega_g
 \leq(2h+1)N^{-1}\sum_{t=1}^Ng_t^2.
\]
The fitted centered post residual differs from \(a_q(Z_{0t}-Z_{jt})\) by
\(-(\widehat\alpha_j-\alpha_j)W_{jt}\) and a sample-centering term.  Therefore
\[
 \widehat\Omega_{v,j}
 =O_P\!\left[
 h_q\left\{a_q^2+\frac{(p_q+a_q)^2}{q}\right\}
 \right].
\]
The fitted bridge sensitivity contains the estimated donor mean and obeys
\[
 \widehat H_j=O_P(q^{-1/2}).
\]
The fitted bridge score before multiplication by \(\widehat H_j\) has conditional second moment \(O_P(p_q^2+a_q^2)\).  Applying the same Bartlett bound gives
\[
 \widehat\Omega_{\psi,j}
 =O_P\!\left\{\frac{h_q(p_q^2+a_q^2)}q\right\}.
\]
Since \(|\mathcal B_q|\asymp q\), \(m_q=q\), and \(h_q=O(q^{1/5})\), the standard-error definition implies
\[
 \widehat{\operatorname{se}}_j^2
 =O_P\!\left\{
 \frac{h_qa_q^2}{q}
 +\frac{h_qp_q^2}{q^2}
 \right\}.
\]
With \(p_q=q^{-2}\) and \(a_q=q^{-4}\),
\[
 \frac{\widehat{\operatorname{se}}_j}{p_q}
 =O_P\!\left[
 \left\{\frac{h_qa_q^2}{qp_q^2}
       +\frac{h_q}{q^2}\right\}^{1/2}
 \right]=o_P(1).
\]
Because \(a_q>0\), the fitted residual vector has a nonzero continuous Gaussian component.  The finite Bartlett kernel matrix is positive definite for \(h_q<m_q\), so the post Bartlett quadratic form, and hence the selected standard error, is strictly positive almost surely.  Combining the last two displays with the numerator expansion proves
\[
 T_q\longrightarrow-\infty
\]
in probability on \(\mathcal N_q\), uniformly over the two possible early-selected labels.  Since \(\Pr(\mathcal N_q)\to1\), the same divergence holds unconditionally and after averaging the conditional law given \(\mathcal F_A\).

Choose a deterministic sequence \(b_q\to\infty\) slowly enough that
\[
 \Pr_{P_q^{(q)}}(T_q>-b_q)\longrightarrow0.
\]
For the conditional Kolmogorov distance \(K_q\), evaluation at the moving point \(-b_q\), followed by Jensen's inequality, gives
\[
 \begin{aligned}
 \mathbb E K_q
 &\geq
 \mathbb E\left|
  \Pr\{T_q\leq-b_q\mid\mathcal F_A\}-\Phi(-b_q)
 \right|\\
 &\geq
 \left|
  \Pr(T_q\leq-b_q)-\Phi(-b_q)
 \right|
 \longrightarrow1.
 \end{aligned}
\]
Every Kolmogorov distance is at most one, so \(\mathbb E K_q\to1\).  This is the maximal possible failure and, in particular, implies the weaker lower bound in the theorem it replaces.
